\begin{pastexam}[label={px:df3-111-ds-04}]{地方三等資訊處理 111 年 第四題}
（申論；20 分）

二元堆積（Binary Heap）是一種優先佇列（Priority Queue），主要用來管理具有優先權順序的資料物件，每個資料物件具有
一個可以界定大小或前後順序的鍵值（Key），我們在此假設鍵值越低的資料物件有越高的優先權。
\begin{itemize}[leftmargin=2.2em, nosep]
\item[㈠] 請完整描述最小堆積（Min\_Heap）的定義與相關的操作功能。（5 分）
\item[㈡] 請說明堆積排序（Heap Sort）的方法並分析其時間複雜度。（5 分）
\item[㈢] 若有兩個二元樹 $T_1$ 及 $T_2$，其節點具有堆積特性且高度分別是 $O(\log n)$ 與 $O(\log m)$，請提供一個方法將此
兩個二元樹結合成為一個節點具有堆積特性的二元樹 $T$，此方法的時間須為 $O(\log n + \log m)$。（10 分）
\end{itemize}
\end{pastexam}

\begin{pastsoln}
（㈡ 屬資結書第~14~章，其餘屬第~8~章。）

\textbf{㈠ 最小堆積的定義}（資結書第~8.1~節）：一棵\textbf{完整二元樹}（complete binary tree：除了最後一層都填滿，
最後一層從左邊連續填），而且每個節點的鍵值都\textbf{不大於}它的子節點的鍵值（堆積性質）。 因此根是最小值，也就是
優先權最高的物件。 因為是完整二元樹，用陣列 $A[1..n]$ 存，節點 $i$ 的子節點是 $2i$、$2i + 1$，父節點是
$\lfloor i/2 \rfloor$，高度 $\lfloor \lg n \rfloor$。

操作（$n$ 是目前的元素個數）：
\begin{itemize}
\item \textbf{取最小值}（Find-Min）：傳回 $A[1]$，$O(1)$。
\item \textbf{插入}（Insert）：放到 $A[n+1]$，再\textbf{往上調整}：比父節點小就和父節點交換，直到不比父節點小或到根，
$O(\log n)$。
\item \textbf{刪除最小值}（Delete-Min / Extract-Min）：取出 $A[1]$，把最後一個 $A[n]$ 搬到根，再\textbf{往下調整}：
和較小的子節點比，比它大就交換，一路往下，$O(\log n)$。
\item \textbf{減少鍵值}（Decrease-Key）：改小後往上調整，$O(\log n)$。
\item \textbf{建堆積}（Build-Heap）：從最後一個內部節點 $\lfloor n/2 \rfloor$ 倒回到 $1$ 逐一往下調整，$O(n)$
（資結書第~8.2~節、定理~8.1）。
\end{itemize}

\textbf{㈡ 堆積排序}（資結書第~14.6~節、虛擬碼~14.13）：由小到大排時用最大堆積。
\begin{enumerate}
\item 把 $A[1..n]$ 建成最大堆積，$O(n)$。
\item 重複 $n - 1$ 次：把根（目前最大）和堆積的最後一個元素交換，堆積大小減一，對新的根往下調整。 第 $k$ 次之後，
$A[n-k+1..n]$ 是最大的 $k$ 個數，已排好。
\end{enumerate}
時間：建堆積 $O(n)$，之後每次往下調整最多走樹高 $\lfloor \lg n \rfloor$ 層，共
$\sum_{k=2}^{n} O(\log k) = O(n \log n)$。 最好、平均、最壞都是 $O(n \log n)$（即使輸入已排序，每次搬到根的仍是小數，
還是要往下走）。 只用常數個額外空間（就地排序），但不穩定。 如果要用本題的最小堆積，同樣的做法會把最小值換到最後，
得到由大到小的順序。

\textbf{㈢ 合併。} $T_1$、$T_2$ 只說有堆積性質（以最小堆積為例），不一定是完整二元樹，所以用指標表示，高度
$h_1 = O(\log n)$、$h_2 = O(\log m)$。 做法：拿一個節點當新的根，把兩棵樹接在它底下，再往下調整。
\begin{enumerate}
\item 若其中一棵是空的，答案就是另一棵。
\item 從 $T_2$ 的根一路往下（有左子節點走左邊，否則走右邊）走到一片葉子 $x$，把 $x$ 從 $T_2$ 拆下來（父節點的指標設
為空；$T_2$ 只有一個節點時 $T_2$ 變空）。 拆掉葉子不影響 $T_2$ 其他節點的堆積性質。 花 $O(h_2)$。
\item 令 $x$ 為新的根，$x.\mathit{left} \gets T_1$，$x.\mathit{right} \gets T_2$。
\item 對 $x$ 往下調整：$x$ 的鍵值比較小的那個子節點大時，和它交換鍵值（或交換節點），再從那個子節點繼續往下，直到
不比子節點大或到達葉子。
\end{enumerate}
\textbf{正確性。} 第 4 步開始時，除了根 $x$ 以外，$x$ 的兩棵子樹各自是堆積；往下調整每次把三者中最小的換到上面，
被換下去的那棵子樹又變成「只有根可能違規」的情形，和資結書虛擬碼~8.1 的向下調整相同，結束時整棵樹滿足堆積性質。
節點集合正好是 $T_1$、$T_2$ 的全部節點。

\textbf{時間。} 第 2 步 $O(h_2)$；新樹的高度是 $\max(h_1, h_2) + 1$，第 4 步每層 $O(1)$，最多 $\max(h_1, h_2) + 1$ 層。
合計 $O(h_1 + h_2) = O(\log n + \log m)$，而且新樹的高度也只比原來多 $1$。 （如果 $T_1$、$T_2$ 是左傾樹，也可以用
沿右路徑合併的方法，同樣是 $O(\log n + \log m)$，見資結書第~8.5~節。）
\end{pastsoln}
